\section*{Erd\H{o}s Problem 1156 --- GPT 6 Astra Ultra partial attempt}


\subsection*{FORMAL RESTATEMENT}
Let $G_n$ have distribution $G(n,1/2)$. Can an absolute constant $C$ bound the number of possible chromatic-number values needed to capture probability tending to one as $n\to\infty$? The problem also asks whether a sufficiently slowly diverging $\omega(n)\to\infty$ satisfies, for every centering function $a(n)$,
\[
\mathbb P\bigl(|\chi(G_n)-a(n)|<\omega(n)\bigr)<\tfrac12
\]
for all sufficiently large $n$.

\subsection*{CONTEXT AND SCOPE}
This replaces the earlier statement-only record with an elementary mathematical attempt. The arguments below establish only their stated partial results; no novelty or full solution is claimed.

This record preserves both questions from the official statement. Its remarks already report results obstructing bounded-width concentration; the stronger eventual anti-concentration question must not be silently replaced by an infinitely-often assertion.

\subsection*{ATTACK PLAN AND WORK}
Expose the independent random blocks consisting of the edges from vertex $j$
to vertices $1,\ldots,j-1$, for $j=1,\ldots,n$. Changing one block changes only
edges incident with one vertex. For the two resulting graphs, deleting that
vertex leaves the same graph, so their chromatic numbers differ by at most one.

This gives a self-contained martingale route to a coarse concentration bound.
The conditional ranges of the corresponding Doob increments have length at most
one, by coupling the unexposed blocks. For a centered variable with range length
one, the exponential bound $\mathbb E e^{tX}\le e^{t^2/8}$ follows from convexity
of the exponential (the elementary bounded-range exponential inequality).
Iterating the conditional inequality and applying Markov's inequality gives
\[\Pr(|\chi(G_n)-\mathbb E\chi(G_n)|\ge t)\le2e^{-2t^2/n}.\tag{1}\]
For example, an interval of radius $\sqrt{n\log n}$ captures probability at least
$1-2n^{-2}$. No independence of vertex colors is being assumed here.

This is an upper concentration estimate on a much larger scale than the question.
It is logically compatible both with bounded-width concentration and with
slowly growing unavoidable fluctuations. It cannot be inverted to obtain the
requested eventual lower bound on the spread.

\subsection*{VERIFICATION AND REMAINING GAP}
Prove anti-concentration at every sufficiently large size, with the quantifiers
over centers and slowly diverging widths in the source statement. The exposure
martingale only bounds large deviations from above.

\subsection*{SOURCES}
https://www.erdosproblems.com/1156

\subsection*{FINAL}
\textbf{UNRESOLVED.} The full source problem remains outside the scope of the proved partial results.

\subsection*{COMPLETION ESTIMATE}
$5\%$. This is a conservative subjective estimate toward the whole problem, not a probability that the argument is correct or a claim that a fixed fraction of the problem has been solved.
